\section{Présentation du projet}

FeastVerse est un projet universitaire réalisé dans le cadre d'un enseignement consacré au développement d'API avec Java Spring Boot. Le sujet imposait essentiellement le framework Spring Boot et la mise en place d'une authentification JWT ; le domaine fonctionnel restait libre.

J'ai choisi d'utiliser cette liberté pour aller nettement au-delà d'un exercice limité à quelques opérations CRUD protégées par un jeton. J'ai imaginé FeastVerse comme le back-end d'une plateforme de recettes intégrant des problématiques proches d'une application réelle : gestion d'utilisateurs et de rôles, propriété et visibilité des ressources, commentaires, likes, signalements et modération, recherche multicritère, pagination, upload d'images, documentation détaillée et déploiement dans un environnement distant.

Cette ambition a progressivement transformé le sujet pédagogique initial en un projet plus large de conception d'API. FeastVerse est mobilisé dans ce dossier comme preuve principale de développement back-end applicatif : il me permet de présenter une solution conçue de bout en bout, de son architecture logicielle jusqu'à sa documentation et sa mise à disposition sur Railway.

\section{Compétences mobilisées}

\begin{table}[H]
\centering
\begin{tabular}{p{3cm} p{11cm}}
\toprule
\textbf{Compétence} & \textbf{Mobilisation dans le projet FeastVerse} \\
\midrule
C1.3 & Définition d'un périmètre fonctionnel dépassant le sujet initial, identification des règles métier et des risques liés à l'augmentation du périmètre. \\
C2.1 & Conception de l'architecture applicative autour de contrôleurs, services, repositories, DTO, mappers et composants de sécurité. \\
C2.2 & Structuration du code, validation des entrées, gestion centralisée des erreurs, versionnement et séparation des responsabilités. \\
C2.4 & Développement back-end Java Spring Boot, API REST, ORM, PostgreSQL, sécurité JWT et logique métier. \\
C3.2 & Validation fonctionnelle à l'aide d'une collection Postman couvrant des cas nominaux, des erreurs et des règles d'autorisation. \\
C3.4 & Déploiement automatique de l'API sur Railway à partir du dépôt GitHub et configuration distincte de l'environnement de production. \\
C3.6 & Documentation technique avec Swagger/OpenAPI et production d'un site MkDocs dédié, lui aussi déployé sur Railway. \\
\bottomrule
\end{tabular}
\caption{Compétences mobilisées dans le projet FeastVerse}
\label{tab:competences-feastverse}
\end{table}

\section{Du sujet pédagogique à un back-end applicatif complet}

Le besoin minimal laissait une grande liberté : développer une API Spring Boot sur un sujet libre et intégrer une authentification JWT. J'aurais donc pu me limiter à quelques entités, des opérations CRUD et un système d'authentification simple.

J'ai volontairement élargi ce périmètre. Le domaine des recettes me permettait de créer des relations métier suffisamment riches pour travailler la sécurité, la recherche, la propriété des ressources et les interactions entre utilisateurs. Le projet a ainsi intégré les utilisateurs et leurs rôles, les ingrédients et leurs types, les recettes et leurs étapes, les commentaires, les likes, les signalements ainsi que plusieurs ressources de référence.

Le périmètre développé comprend également plus de 700 ingrédients de démonstration, trois rôles distincts --- utilisateur, modérateur et administrateur --- ainsi que des mécanismes de pagination, tri et filtrage sur les principales ressources. Les utilisateurs peuvent interagir avec les contenus, tandis que les modérateurs et administrateurs disposent de droits spécifiques sur les signalements et certaines suppressions.

Cette ambition a toutefois eu un coût. Toutes les ressources n'ont pas atteint le même niveau de complétude dans le temps disponible : les utilisateurs, ingrédients, commentaires, likes et signalements sont les plus avancés, tandis que les recettes et leurs étapes étaient principalement disponibles en lecture au moment du rendu initial. Je préfère présenter cette limite explicitement plutôt que qualifier FeastVerse d'application totalement finalisée.

\section{Conception de l'architecture back-end}

\subsection{Stack et organisation générale}

Le cadre pédagogique imposait Spring Boot ; mon travail a donc surtout porté sur la manière de structurer l'application autour de ce framework. J'ai utilisé Gradle pour le build, Spring Web pour les endpoints REST, Spring Data JPA et Hibernate pour la persistance, Spring Security pour les autorisations, JWT pour l'authentification, PostgreSQL pour la base de données et MapStruct pour les conversions entre entités et DTO.

J'ai organisé l'application en couches afin d'éviter que les contrôleurs concentrent à la fois la gestion HTTP, la logique métier et l'accès aux données. Les contrôleurs reçoivent les requêtes et délèguent aux services ; les services portent les règles métier ; les repositories restent centrés sur la persistance. Les DTO et mappers séparent quant à eux le modèle de stockage des contrats exposés par l'API.

\begin{figure}[H]
\centering
\fbox{\parbox{0.88\textwidth}{\centering Emplacement prévu : schéma de l'architecture FeastVerse.\\Client HTTP / Postman $\rightarrow$ contrôleurs REST $\rightarrow$ services $\rightarrow$ repositories $\rightarrow$ PostgreSQL.\\DTO et mappers entre les couches d'exposition et le modèle de persistance.}}
\caption{Architecture générale de l'API FeastVerse}
\label{fig:architecture-feastverse}
\end{figure}

Cette séparation s'est révélée particulièrement utile pour les utilisateurs. Selon le contexte, une même ressource ne doit pas exposer les mêmes informations : une vue publique, une vue personnelle et une vue administrateur peuvent donc utiliser des DTO différents sans modifier l'entité de persistance.

\subsection{Persistance et modèle métier}

Le modèle relationnel traduit les interactions entre utilisateurs et contenus : un utilisateur peut créer certaines ressources, aimer une recette, commenter un contenu ou créer un signalement. Une recette est liée à des ingrédients, des étapes, des commentaires, des likes et des signalements.

Ces relations m'ont obligé à traiter des règles qui dépassent un CRUD simple : vérifier le propriétaire d'une ressource, distinguer visibilité publique et privée, gérer la suppression logique ou définitive et appliquer des autorisations différentes selon le rôle connecté.

\section{Concevoir une API interrogeable à grande échelle}

Une partie importante du projet concerne la consultation des ressources. Avec plusieurs centaines d'ingrédients et de nombreux critères métier, renvoyer des listes complètes n'était pas une solution satisfaisante. J'ai donc généralisé pagination, tri et filtrage sur les principales routes de consultation.

L'exemple des recettes est le plus représentatif. Le contrôleur accepte plusieurs critères optionnels : nom, temps maximal, difficulté, type, ingrédients inclus ou exclus, tags, propriétaire et visibilité. L'objet \texttt{Pageable} permet de faire varier pagination et tri sans créer un endpoint par combinaison de critères.

\paragraph{Extrait de code : filtrage et pagination des recettes}

\begin{lstlisting}[language=Java,caption={Endpoint de recherche paginée et filtrée des recettes},label={lst:feastverse-filtrage-recettes}]
@GetMapping
public Page<RecipeViewDto> getAllRecipes(
        @RequestParam(required = false) String name,
        @RequestParam(required = false) Integer maxTotalTime,
        @RequestParam(required = false) RecipeDifficulty difficulty,
        @RequestParam(required = false) UUID typeId,
        @RequestParam(required = false) List<Long> withIngredient,
        @RequestParam(required = false) List<Long> withIngredientType,
        @RequestParam(required = false) List<Long> withoutIngredientType,
        @RequestParam(required = false) List<String> withTags,
        @RequestParam(required = false) UUID ownerId,
        @RequestParam(required = false, defaultValue = "PUBLIC")
        VisibilityFilter visibility,
        @ParameterObject
        @PageableDefault(size = 10, sort = "likeCount",
                         direction = Sort.Direction.DESC)
        Pageable pageable
) {
    Page<Recipe> recipes = recipeService.findAllRecipes(
            name, maxTotalTime, difficulty, typeId,
            withIngredient, withIngredientType,
            withoutIngredientType, withTags,
            ownerId, visibility, pageable
    );
    return recipes.map(recipeMapper::toRecipeListViewDto);
}
\end{lstlisting}

Ce choix permet de combiner plusieurs usages dans un même contrat HTTP tout en limitant le volume transféré. La conversion vers un DTO de liste évite également de renvoyer l'entité complète lorsque seules quelques informations sont nécessaires à l'affichage d'un résultat.

\section{Sécurité et autorisations métier}

\subsection{Authentification JWT}

L'authentification JWT faisait partie des contraintes du sujet. J'ai cependant poussé son utilisation au-delà de la seule génération d'un jeton. Un filtre dédié extrait le token transmis dans l'en-tête \texttt{Authorization}, vérifie sa validité, recharge l'utilisateur puis initialise le contexte de sécurité Spring avant l'exécution du contrôleur.

\begin{lstlisting}[language=Java,caption={Extraction et validation du JWT dans le filtre de sécurité},label={lst:feastverse-jwt-filter}]
final String authorizationHeader = request.getHeader("Authorization");

String jwt = null;
UUID userId = null;

if (authorizationHeader != null && authorizationHeader.startsWith("Bearer ")) {
    jwt = authorizationHeader.substring(7);
    userId = jwtUtil.extractUserId(jwt);
}

if (userId != null &&
    SecurityContextHolder.getContext().getAuthentication() == null) {

    UserDetails userDetails = userDetailsService.loadUserById(userId);

    if (jwtUtil.validateToken(jwt, userId)) {
        UsernamePasswordAuthenticationToken authToken =
                new UsernamePasswordAuthenticationToken(
                        userDetails,
                        null,
                        userDetails.getAuthorities()
                );
        SecurityContextHolder.getContext().setAuthentication(authToken);
    }
}
\end{lstlisting}

\subsection{Rôles, propriété et modération}

J'ai défini trois niveaux principaux d'autorisation : utilisateur standard, modérateur et administrateur. Ils ne donnent pas accès aux mêmes opérations. Les utilisateurs peuvent créer certaines ressources et interagir avec les contenus ; les modérateurs disposent d'actions liées au traitement des signalements ; les administrateurs peuvent effectuer des opérations plus sensibles, notamment certaines suppressions définitives.

Les annotations \texttt{@PreAuthorize} rendent une partie de ces règles visible au niveau des endpoints, tandis que d'autres contrôles tiennent compte du propriétaire de la ressource. L'objectif était d'obtenir un modèle où être authentifié ne signifie pas automatiquement avoir le droit d'agir sur n'importe quelle donnée.

\begin{lstlisting}[language=Java,caption={Exemples de restrictions d'accès par rôle},label={lst:feastverse-preauthorize}]
@PreAuthorize("hasRole('ROLE_STANDARD')")
@PostMapping(consumes = MediaType.MULTIPART_FORM_DATA_VALUE)
public IngredientViewDto createIngredient(...) { ... }

@PreAuthorize("hasAnyRole('ROLE_MODERATOR', 'ROLE_ADMINISTRATOR')")
@GetMapping("/reports")
public Page<ReportViewDTO> getAll(...) { ... }

@PreAuthorize("hasRole('ROLE_ADMINISTRATOR')")
@DeleteMapping("/reports/{reportId}")
public ResponseEntity<Void> deleteReport(...) { ... }
\end{lstlisting}

La gestion des ingrédients illustre aussi cette logique : création avec upload d'image, visibilité et propriété de la ressource, suppression logique pour préserver les données, et suppression définitive réservée à l'administrateur.

\section{Qualité des réponses et validation fonctionnelle}

\subsection{Gestion centralisée des erreurs}

J'ai centralisé le traitement de plusieurs erreurs afin que les clients reçoivent des réponses cohérentes plutôt que des exceptions techniques brutes. Les erreurs de validation, ressources inexistantes, formats invalides ou conflits métier sont transformés en réponses HTTP structurées.

\begin{lstlisting}[language=Java,caption={Gestion centralisée des erreurs de validation},label={lst:feastverse-validation-handler}]
@ExceptionHandler(MethodArgumentNotValidException.class)
public ResponseEntity<ErrorResponse> handleValidationExceptions(
        MethodArgumentNotValidException ex,
        WebRequest request
) {
    Map<String, String> errors = new HashMap<>();

    for (FieldError error : ex.getBindingResult().getFieldErrors()) {
        errors.put(error.getField(), error.getDefaultMessage());
    }

    ErrorResponse response = new ErrorResponse(
            HttpStatus.BAD_REQUEST.value(),
            HttpStatus.BAD_REQUEST.getReasonPhrase(),
            "Validation failed",
            getPath(request),
            errors
    );

    return ResponseEntity.status(HttpStatus.BAD_REQUEST).body(response);
}
\end{lstlisting}

\subsection{Scénarios Postman}

La validation fonctionnelle s'est principalement appuyée sur une collection Postman organisée par domaine. J'y ai construit des scénarios couvrant authentification, rôles, consultation et modification de ressources, filtrage, pagination, upload multipart, suppression logique et définitive, likes, commentaires et signalements.

Certaines requêtes réutilisent les résultats d'étapes précédentes. Après authentification, par exemple, un script vérifie la réponse et stocke automatiquement le JWT dans une variable de collection pour les appels suivants.

\begin{lstlisting}[language=JavaScript,caption={Script Postman de vérification et de stockage du token},label={lst:feastverse-postman-token}]
const jsonData = pm.response.json();

pm.test("Status code is 200", function () {
    pm.expect(pm.response.code).to.equal(200);
});

pm.test("Response contains a user-token", function () {
    pm.expect(jsonData).to.have.property("accessToken").that.is.not.empty;
});

if (jsonData.accessToken) {
    pm.collectionVariables.set("user-token", jsonData.accessToken);
}
\end{lstlisting}

Cette collection constitue un plan de validation fonctionnelle reproductible, mais elle ne remplace pas une véritable suite automatisée de tests unitaires et d'intégration. C'est l'une des limites que je retiens du projet.

\section{Documenter pour rendre l'API utilisable}

\subsection{Swagger et OpenAPI}

Je ne voulais pas que l'utilisation de l'API dépende de la lecture du code source. J'ai donc documenté les contrôleurs avec OpenAPI en précisant les paramètres, réponses, erreurs et exemples de payloads.

Certaines réponses variant selon le rôle de l'utilisateur, j'ai créé des annotations de documentation dédiées. Pour la consultation d'un utilisateur, par exemple, la documentation distingue les schémas public, personnel et administrateur et fournit des exemples de réponse correspondants. Cette démarche a été appliquée progressivement aux différents contrôleurs de l'application.

\begin{figure}[H]
\centering
\fbox{\parbox{0.88\textwidth}{\centering Emplacement prévu : capture de Swagger UI montrant un endpoint documenté avec paramètres, réponses et exemples.}}
\caption{Documentation Swagger/OpenAPI de l'API FeastVerse}
\label{fig:swagger-feastverse}
\end{figure}

\subsection{Site MkDocs dédié}

J'ai également produit un site de documentation séparé avec MkDocs. Il regroupe une documentation d'authentification et des endpoints, des pages fonctionnelles et cas d'usage, une roadmap, un changelog, une présentation de l'architecture, de la stack technique ainsi qu'un glossaire.

Cette documentation répond à un besoin différent de Swagger. OpenAPI décrit précisément les contrats HTTP ; MkDocs permet de présenter le projet, son organisation et son usage avec davantage de recul. Le site intègre notamment Mermaid et un composant Swagger UI.

J'ai ensuite conteneurisé cette documentation avec un build multi-stage : une première image Python génère le site MkDocs en mode strict, puis une image Nginx légère sert les fichiers statiques produits. La documentation a elle aussi été déployée sur Railway.

\begin{figure}[H]
\centering
\fbox{\parbox{0.88\textwidth}{\centering Emplacement prévu : capture de la documentation MkDocs déployée, idéalement la page Architecture ou la page d'accueil.}}
\caption{Documentation MkDocs déployée du projet FeastVerse}
\label{fig:mkdocs-feastverse}
\end{figure}

\section{Déploiement automatique sur Railway}

J'ai déployé l'API sur Railway et relié le service à son dépôt GitHub afin qu'une nouvelle version soit reconstruite et redéployée automatiquement après mise à jour de la branche suivie. Cette étape m'a obligé à distinguer les paramètres locaux des paramètres de production et à externaliser les éléments sensibles.

Le profil de production utilise des variables d'environnement pour la connexion PostgreSQL, le secret JWT et les paramètres Cloudinary. Le service peut ainsi être reconstruit dans Railway sans intégrer les secrets au dépôt.

\begin{lstlisting}[language=yaml,caption={Extrait de configuration de production},label={lst:feastverse-config-prod}]
spring:
  datasource:
    url: jdbc:postgresql://${PGHOST}:${PGPORT}/${POSTGRES_DB}
    username: ${POSTGRES_USER}
    password: ${POSTGRES_PASSWORD}

jwt:
  secret: ${JWT_SECRET}

upload:
  storage-type: cloudinary
  cloudinary:
    cloud-name: ${CLOUDINARY_CLOUD_NAME}
    api-key: ${CLOUDINARY_API_KEY}
    api-secret: ${CLOUDINARY_API_SECRET}
\end{lstlisting}

Ce déploiement ne constitue pas à lui seul une chaîne DevOps complète, mais il apporte une preuve concrète de déploiement continu : le code versionné reste la source du service, Railway automatise la reconstruction et la mise en ligne, et la configuration de production est séparée du code applicatif.

\begin{figure}[H]
\centering
\fbox{\parbox{0.88\textwidth}{\centering Emplacement prévu : historique Railway montrant plusieurs déploiements associés aux commits GitHub, sans afficher de variable sensible.}}
\caption{Déploiements automatiques de FeastVerse sur Railway}
\label{fig:railway-feastverse}
\end{figure}

\section{Limites et bilan critique}

Le principal enseignement du projet tient autant à son ambition qu'à ses limites. En partant d'un exercice Spring Boot avec JWT, j'ai choisi d'explorer un périmètre beaucoup plus large : architecture en couches, autorisations métier, ressources possédées, modération, recherche avancée, upload d'images, documentation OpenAPI, site MkDocs et déploiement automatique sur Railway.

Cette ambition m'a permis d'aborder des problématiques beaucoup plus riches qu'un CRUD classique, mais elle a aussi dispersé le temps disponible. Le CRUD complet des recettes et des étapes, pourtant central dans le domaine choisi, n'a pas été terminé au même niveau que les autres ressources. De même, les scénarios Postman sont nombreux, mais la couverture automatisée par tests unitaires et d'intégration reste insuffisante.

Avec le recul, je conserverais cette ambition technique, mais je définirais plus tôt un socle fonctionnel minimal à terminer intégralement avant d'ajouter les fonctionnalités secondaires. Le CRUD complet des recettes, les règles d'autorisation essentielles et une première suite de tests automatisés constitueraient les critères de sortie d'un premier incrément ; la modération, l'upload, la documentation enrichie et les fonctionnalités complémentaires viendraient ensuite.

FeastVerse m'a ainsi permis de consolider mes compétences en conception et développement back-end, mais surtout de comprendre qu'aller plus loin techniquement n'est pertinent que si le périmètre reste maîtrisé. Le projet constitue une preuve forte de développement d'API, de sécurisation, de documentation et de mise à disposition automatisée, tout en offrant un recul concret sur la priorisation et la notion de produit minimum viable.
